% 无阻塞潜能语义
\section{无阻塞潜能：语义与边界}\label{sec:bg-nonblocking}

本文判定的「无阻塞」是\textbf{无阻塞潜能}，而非保证：

\begin{quote}
在最优自适应路由下，物理资源足以支撑所有指定流量模式的同时交换——网络\textbf{有可能}达到无阻塞。
\end{quote}

这里有两层放松。其一，路由是最优的：LP 为每个流量模式选择最佳分流方案，
若最优路由都不够，任何次优路由更不够；
其二，负载是\textbf{包络}而非等式：物理资源按最坏模式配置，路由可灵活利用未被某模式用满的链路。
因此这是早期设计的\textbf{早筛}判据——假阴性（该过的没过）不可接受，假阳性（该剪的没剪）留给后续 NoC 仿真消化。

与之相对的是\textbf{可重排无阻塞（RNB）}：它对任意连接请求给出调度无关的硬保证，条件更强（充分而非必要），
且不依赖拓扑对称性。本文以「潜能」为主，RNB 可作为可行域下界在实验中对比。

需要说明方法边界：本文的潜能判定预设物理拓扑\textbf{高度对称}（vertex-transitive），
从而把 $n!$ 个排列压缩到可控数量的共轭轨道代表元（群论归约的具体推导见附录 A）。
对称性不足的拓扑不在本方法讨论范围内；这是以对称性假设换取线性规划可解性的取舍，
在此边界内方法是严格的。
